\documentclass[11pt]{article}

\usepackage[T1]{fontenc}
\usepackage{lmodern}
\usepackage{amsmath,amssymb,amsthm}
\usepackage{booktabs}
\usepackage{enumitem}
\usepackage[hidelinks]{hyperref}
\usepackage[margin=1in]{geometry}

\hypersetup{
  pdftitle={Tuza's Conjecture for Graphs with at Most Four Edge-Disjoint Triangles},
  pdfauthor={Samuel Masse},
  pdfsubject={A bounded-packing case of Tuza's triangle packing-covering conjecture},
  pdfkeywords={Tuza conjecture, triangle packing, triangle transversal, SAT, LRAT, Lean}
}

\newtheorem{theorem}{Theorem}[section]
\newtheorem{lemma}[theorem]{Lemma}
\newtheorem{proposition}[theorem]{Proposition}
\newtheorem{corollary}[theorem]{Corollary}
\newtheorem{claim}[theorem]{Claim}
\theoremstyle{definition}
\newtheorem{definition}[theorem]{Definition}
\newtheorem{remark}[theorem]{Remark}

\newcommand{\Tcal}{\mathcal{T}}
\newcommand{\supp}{\operatorname{supp}_{\triangle}}
\newcommand{\defect}{\operatorname{def}}

\title{Tuza's Conjecture for Graphs with at Most Four\\Edge-Disjoint Triangles}
\author{Samuel Masse}
\date{}

\begin{document}

\maketitle
\vspace{-2em}

\begin{abstract}
For a finite simple graph $G$, let $\nu(G)$ be the maximum number of
pairwise edge-disjoint triangles and let $\tau(G)$ be the minimum number of
edges meeting every triangle.  We prove that $\tau(G)\leq 2\nu(G)$ whenever
$\nu(G)\leq 4$.  The cases $\nu(G)\leq 3$ were proved by Parker.  For
$\nu(G)=4$, a discharging reduction based on the method of Sheng and Xiao
shows that a counterexample would have triangle support on at most twelve
vertices.  Marking a maximum four-packing leaves 66 finite order-type cases.
We encode these cases as SAT instances and certify each resulting formula by
an LRAT refutation.  A Lean formalization reconstructs the semantic clauses,
checks the exact formulas and all 66 refutations, and derives the stated
graph theorem.
\end{abstract}

\section{Introduction}

A triangle packing in a graph is a family of pairwise edge-disjoint
triangles, and a triangle transversal is a set of graph edges meeting every
triangle.  Tuza conjectured that every finite simple graph $G$ satisfies
\begin{equation}\label{eq:tuza}
  \tau(G)\leq 2\nu(G).
\end{equation}
The factor two is best possible already for $K_4$ and $K_5$.  The conjecture
remains open in general; see Tuza's original paper~\cite{Tuza1990} and the
subsequent literature cited by Parker~\cite{Parker2025}.

The best universal quantitative bound currently available is weaker than
factor two; Haxell's bound~\cite{Haxell1999}, for example, gives only
$\tau(G)\leq11$ when $\nu(G)=4$.  Many exact results instead impose a graph
class, such as maximum average degree below seven~\cite{Puleo2015} or maximum
degree at most seven~\cite{Gupta2026}.  Parker proved the exact
bounded-packing statement
\begin{equation}\label{eq:parker}
  \nu(G)\leq 3 \quad\Longrightarrow\quad \tau(G)\leq 2\nu(G).
\end{equation}
We establish the next bounded-packing case.

\begin{theorem}\label{thm:main}
Every finite simple graph $G$ with $\nu(G)\leq 4$ satisfies
$\tau(G)\leq 2\nu(G)$.
\end{theorem}

The proof has an analytic part and a computer-assisted finite part.  The analytic
part adapts the $3k$ discharging framework of Sheng and
Xiao~\cite{ShengXiao2024}.  Starting from a maximal packing saturated under
three local exchanges, the argument either finds an exact fat-head crown or
bounds the triangle support by three times the packing size.  Exact offset
identities for crowns and triangle-isolated copies of $K_4$ then give the
following reduction.

\begin{proposition}\label{prop:twelve}
If $\nu(G)=4$ and $\tau(G)\geq 9$, then the triangle support of $G$ has at
most twelve vertices.
\end{proposition}

For the finite part, mark four edge-disjoint triangles.  Their vertex
intersection pattern is a linear $3$-uniform hypergraph with four edges.
There are 66 order-type scopes on orders six through twelve.  In each scope,
all other graph edges remain free.  A SAT formula represents the literal graph
and its complete triangle set.  Whenever a model has five edge-disjoint
triangles or a transversal of at most eight edges, a semantic clause recording
that exact witness is added.  Exhaustion yields 66 final formulas.  Their DRAT
proofs were converted to LRAT and checked in Lean; the original DRAT traces
were also checked independently with
\texttt{drat-trim}~\cite{HeuleWetzlerHunt2014}.

The finite verification does not rely on a catalogue of unlabeled graphs or
on the solver's UNSAT declaration.  The reduction to twelve vertices,
completeness of the marked support types, graph-to-CNF translation, and case
dispatch are part of the Lean proof.

\section{Preliminaries}

All graphs are finite, undirected, loopless, and simple.  Write $\Tcal(G)$
for the set of triangles of $G$.  We identify a triangle with either its
three vertices or its three graph edges when no ambiguity is possible.

The \emph{triangle packing number} $\nu(G)$ is the largest cardinality of a
subfamily of $\Tcal(G)$ whose members are pairwise edge-disjoint.  The
\emph{triangle transversal number} $\tau(G)$ is the minimum size of a set
$C\subseteq E(G)$ such that $C$ meets the edge set of every triangle.
Define the integral Tuza defect
\[
  \defect(G)=\tau(G)-2\nu(G).
\]
A strict counterexample to~\eqref{eq:tuza} is therefore a graph with
$\defect(G)>0$.

The \emph{triangle support} $\supp(G)$ is the graph whose edge set consists
of the edges lying in at least one triangle of $G$, and whose vertex set
consists of their endpoints.  Then
\begin{equation}\label{eq:support}
  \nu(\supp(G))=\nu(G),\qquad
  \tau(\supp(G))=\tau(G).
\end{equation}

For a vertex $v$, form the auxiliary graph whose vertices are the graph edges
incident with $v$, with two such edges adjacent when they lie together in a
triangle.  If this auxiliary graph is disconnected, split $v$ into one vertex
for each component, attaching each incident edge to the corresponding copy.
The edge ground set is unchanged, and three edges form a triangle before the
split exactly when their corresponding edges form a triangle afterward.
Consequently a sequence of such splits preserves the entire triangle
hypergraph, and hence preserves both $\nu$ and $\tau$.  We call the resulting
graph \emph{split-normal}.

We need one further reduction object from Sheng and Xiao.

\begin{definition}[fat-head crown]\label{def:crown}
A fat-head crown in a graph $G=(V,E)$ is a triple $(C,H,X)$ such that $C$ and
$X$ partition $V$, $H\subseteq E$, and:
\begin{enumerate}[label=(\roman*)]
  \item $C$ is independent;
  \item $H$ is exactly the set of edges spanned by vertices of $C$, meaning
        $uv\in H$ exactly when $cuv$ is a triangle for some $c\in C$;
  \item no vertex of $C$ is adjacent to a vertex of $X\setminus V(H)$; and
  \item there is an edge-disjoint packing $P$ of $|H|$ triangles, each with
        exactly one vertex in $C$ and exactly one edge in $H$.
\end{enumerate}
The packing $P$ is a \emph{witness packing}.
\end{definition}

The exact witness requirement is essential.  A subset passing a relaxed
matching test need not be a crown.

\begin{lemma}[exact crown offsets]\label{lem:crown-offset}
Let $(C,H,X)$ be a fat-head crown with witness packing $P$, put $h=|H|$, and
let $G'=G-C-H$.  Then
\[
  \nu(G)=\nu(G')+h,
  \qquad
  \tau(G)=\tau(G')+h.
\]
\end{lemma}

\begin{proof}
Adjoining $P$ lifts every packing of $G'$.  Conversely, delete from a packing
of $G$ every triangle meeting $C$ or $H$.  Each deleted triangle contains a
head edge, and distinct packed triangles contain distinct head edges, so at
most $h$ triangles are deleted.

Adjoining $H$ lifts every transversal of $G'$.  Conversely, a transversal of
$G$ meets the $h$ pairwise edge-disjoint triangles of $P$ in at least $h$
distinct graph edges.  These edges disappear on intersection with $E(G')$,
and that intersection remains a transversal of $G'$.  Both identities follow.
\end{proof}

\begin{lemma}[triangle-isolated $K_4$ edge block]\label{lem:k4-offset}
Suppose a copy $Q\cong K_4$ has no edge in a triangle outside the four faces
of $Q$, and let $G'=G-E(Q)$.  Then
\[
  \nu(G)=\nu(G')+1,
  \qquad
  \tau(G)=\tau(G')+2.
\]
\end{lemma}

\begin{proof}
At most one face of $Q$ can occur in an edge-disjoint packing, while any one
face can be added to a packing of $G'$.  A transversal of the four faces uses
at least two local edges, and any pair of opposite edges meets all four faces.
The isolation hypothesis separates these local choices from $G'$.
\end{proof}

\section{The twelve-vertex reduction}\label{sec:reduction}

We isolate the form of the Sheng--Xiao discharging argument needed here.  Let
$S$ be a maximal triangle packing, let $F=V(G)\setminus V(S)$, and let
$R=E(G)\setminus E(S)$.  We allow the following packing-only improvements:
\begin{enumerate}[label=(E\arabic*)]
  \item replace one member of $S$ by at least two triangles using only $R$
        and the three edges of the removed triangle;
  \item replace two members of $S$ by at least three triangles using only
        $R$ and the six edges of the removed triangles;
  \item replace two members $T_1,T_2$ by two triangles $T'_1,T'_2$ such that
        $(S\setminus\{T_1,T_2\})\cup\{T'_1,T'_2\}$ is a packing whose vertex
        support is strictly larger than the vertex support of $S$.
\end{enumerate}
After each exchange, extend greedily to a maximal packing.  The lexicographic
pair $(|S|,|V(S)|)$ strictly increases and is bounded by
$(\nu(G),3\nu(G))$, so the process terminates.

An edge $uv\in E(S)$ is \emph{labelled} if some $x\in F$ spans it, that is,
$xuv$ is a triangle.  Let $L$ be the set of labelled edges.

\begin{lemma}[discharge bound]\label{lem:discharge}
Let $G$ be triangle-supported and split-normal.  Assume that:
\begin{enumerate}[label=(\alph*)]
  \item $G$ has no triangle-isolated $K_4$ edge block as in
        Lemma~\ref{lem:k4-offset};
  \item $S$ is a maximal packing saturated under (E1)--(E3); and
  \item there is no nonempty fat-head crown $(C,H,X)$ with
        $C\subseteq V(G)\setminus V(L)$ and $H\subseteq L$.
\end{enumerate}
Then $|V(G)|\leq 3|S|$.
\end{lemma}

\begin{proof}
We give the charge argument in full.

First, each triangle of $S$ has at most one labelled edge.  If two of its
edges are spanned by distinct vertices of $F$, (E1) applies.  Suppose instead
that $x\in F$ spans $uv$ and $vw$ in the packed triangle $uvw$.  The vertices
$u,v,w,x$ induce $K_4$.  Since it is not a triangle-isolated edge block, some
edge belongs to an external triangle.  Up to interchanging $u$ and $w$, there
are four edge orbits, and all four give a permitted exchange:
\begin{enumerate}[label=(\roman*)]
  \item if the labelled edge $vw$ lies in $vwy$, use $xuv$ and $vwy$ for
        (E1) when neither $vy$ nor $wy$ lies in $E(S)$.  If $vy$ lies in a
        retained packing triangle, retain that triangle and insert $xuw$;
        if $wy$ lies in a retained packing triangle, retain it and insert
        $xuv$.  Either choice is a safe (E3) exchange, even when both edges
        belong to retained packing triangles;
  \item if the outside spoke $xw$ lies in $xwy$, maximality forces $wy$ into
        a packing triangle; retain that triangle and insert $xuv$ by (E3);
  \item if the opposite edge $uw$ lies in $uwy$, use $xuv$ and $uwy$ for
        (E1), unless $uy$ or $wy$ belongs to $E(S)$; in those cases retain
        its packing triangle and insert $xvw$ or $xuv$, respectively;
  \item if the central spoke $xv$ lies in $xvy$, maximality forces $vy$ into
        a packing triangle; retain it and insert $xuw$.
\end{enumerate}
Here an external triangle means a triangle other than the four faces of the
displayed $K_4$.  In every two-for-two move, the inserted face contains one
edge of a removed packing triangle and two spokes incident with $x\in F$.
It is therefore
edge-disjoint from every retained packing triangle.  Moreover the retained
triangle keeps every old vertex omitted by the inserted face, while the new
face adds $x\notin V(S)$.  Each two-for-two move is consequently a globally
safe (E3) exchange with strictly larger vertex union.  Thus at most one edge
of each member of $S$ is labelled.

Call a packing triangle \emph{bad} if it has no labelled edge and
\emph{good} otherwise.  Delete $L$ temporarily.  A good triangle is
\emph{excellent} if neither unlabelled edge lies in another remaining
triangle, and \emph{pretty-good} otherwise.  Write $S_1,S_2,S_3$ for the
excellent, pretty-good, and bad triangles.  Put
\[
  V_1=V(S_1)\setminus\bigl(V(L)\cup V(S_2)\cup V(S_3)\bigr),
  \qquad
  V_2=V(S_2)\setminus V(L).
\]

The third vertex of every good packing triangle lies outside $V(L)$.  Indeed,
write such a triangle as $uvw$, with labelled edge $uv$ spanned by
$q\in F$, and suppose that $w\in V(L)$.  A different retained packing
triangle then contains a labelled edge incident with $w$.  Retain that
triangle and replace $uvw$ by $quv$.  The result is a packing: $uv$ belonged
to the removed triangle, while the two edges incident with $q\notin V(S)$
belong to no member of $S$.  The retained triangle preserves $w$, the new
triangle preserves $u,v$, and $q$ is a new support vertex.  This is an (E3)
improvement, a contradiction.  In particular, every component meeting
$S_2$ contains a vertex of $V_2$.

Place three units on each labelled edge and on each bad triangle.  Since every
member of $S_1\cup S_2$ has exactly one labelled edge, the total charge is
exactly $3|S|$.  Move one unit from each labelled edge to each endpoint, and
move one unit from each bad triangle to each of its vertices.  Each labelled
edge retains one unit.  At this point only vertices in $F\cup V_1\cup V_2$
can have charge zero.

Consider a connected component of the graph formed by the edges of $S$.  If
it contains $x$ packing triangles and covers $m$ vertices, form the bipartite
incidence graph between those triangles and those vertices.  This incidence
graph is connected, has $x+m$ vertices, and has exactly $3x$ edges.  Hence
$x+m\leq 3x+1$, so $m\leq 2x+1$.  If $x_2$ of the component's triangles are
bad, the charge available on its vertices after the first move is $2x+x_2$.
Move surplus charge inside the component, prioritizing vertices of $V_1$.
Each component then has at most one zero; a component meeting $S_3$ has none,
and a component meeting $S_2$ can have a zero only in $V_2$.

We next establish the component transfer and the residual-zero statements
that complete the discharging step; compare
Sheng--Xiao~\cite[Lemmas~8--12]{ShengXiao2024}.

\begin{claim}[componentwise transfer]\label{clm:transfer}
For every packing component whose unique zero is a vertex $w\in V_2$, there
are two labelled edges in that component with one common, unique spanner
$q\in F$.  Moving one of their remaining units to $w$ and the other to $q$
makes $w$ positive.  This can be done once in every such component without
using any labelled edge twice.  Whenever a labelled edge is reduced to zero,
its unique spanner in $F$ receives a unit in the same component transfer.
\end{claim}

\begin{proof}[Proof of the claim]
Choose the pretty-good triangle $wuv\in S_2$ containing $w$, with labelled
edge $uv$ spanned by some $q\in F$.  Since it is pretty-good, after possibly
interchanging $u$ and $v$, the unlabelled edge $uw$ lies in a triangle $uwx$
of $G-L$.  If neither $ux$ nor $wx$ lies in $E(S)$, then $quv$ and $uwx$
replace $wuv$ by (E1), a contradiction.  If $wx$ lies in a packing triangle,
retain that triangle and replace $wuv$ by $quv$.  The two spokes at $q$ avoid
every retained member, and the new pair gains the vertex $q$, so this is a
globally safe (E3), again a contradiction.  We may therefore assume that
$ux$ lies in a packing triangle $uxy$ and that $wx\notin E(S)$.

The triangle $uxy$ is not bad, since its packing component still has a zero
after the component redistribution.  Its labelled edge is not $ux$, because
$uwx$ lies in $G-L$.  If its labelled edge were $xy$, with spanner $q'\in F$,
then retaining $wuv$ and replacing $uxy$ by $q'xy$ would be a globally safe
(E3) exchange whose vertex union gains $q'$.  Hence $uy$ is labelled.  The
triangle $uxy$ is pretty-good because its unlabelled edge $ux$ lies in
$uwx\subseteq G-L$.

The labelled edges $uv$ and $uy$ have one common unique spanner.  Indeed, if
$q_1\in F$ spans $uv$ and $q_2\in F$ spans $uy$ with $q_1\ne q_2$, then the
three triangles
\[
  q_1uv,\qquad q_2uy,\qquad uwx
\]
are pairwise edge-disjoint and use only edges outside $E(S)$ together with
edges of the two removed triangles $wuv$ and $uxy$; recall that
$wx\notin E(S)$.  They give (E2).  Applying this observation with every
choice of spanner shows both uniqueness and equality; call the common vertex
$q$.  Transfer the unit on $uv$ to $w$ and the unit on $uy$ to $q$.

There is at most one zero in a packing component.  We make the displayed
choice only once for that component.  Distinct components contain disjoint
packing edges, so their two labelled edges are disjoint and no charge is
spent twice.  The last assertion follows from the construction: every edge
that loses its last unit is one of the displayed pair, and the common spanner
$q$ receives the other edge's unit in the same transfer.
\end{proof}

Perform the transfer of Claim~\ref{clm:transfer} in every eligible component.
Every vertex in $V_2$ is now positive.  Let $I$ be the zero-charge vertices
remaining in $F\cup V_1$, and let $L'$ be the labelled edges that still carry
one unit.

\begin{claim}[residual independence]\label{clm:independent}
The set $F\cup V_1$, and hence $I$, is independent.
\end{claim}

\begin{proof}[Proof of the claim]
Suppose that $uv$ is an edge with $u,v\in F\cup V_1$.  Triangle support gives
a triangle $uvw$, and maximality of $S$ forces at least one of its three
edges into $E(S)$.  If $uv\in E(S)$, then $u,v\in V_1$ and the packing
triangle containing $uv$ is neither bad nor good: it cannot be bad because
$V_1$ avoids $V(S_3)$, while each possible labelled edge of a good triangle
has endpoint $u$ or $v$, contrary to $V_1\cap V(L)=\varnothing$.

Otherwise, say $uw\in E(S)$.  Then $u\in V_1$, and the packing triangle
containing $uw$ belongs to $S_1$, since $u$ lies in no triangle of
$S_2\cup S_3$.  None of $uv,vw,uw$ is labelled because neither $u$ nor $v$
lies in $V(L)$.  Thus the unlabelled packing edge $uw$ lies in the other
triangle $uvw$ of $G-L$, contradicting excellence.  Both cases are
impossible.
\end{proof}

\begin{claim}[outside residual zeros]\label{clm:fzero}
Every vertex of $I\cap F$ spans only edges of $L'$.
\end{claim}

\begin{proof}[Proof of the claim]
A vertex of $F$ can span only an edge of $E(S)$: otherwise its triangle would
be edge-disjoint from the maximal packing.  Such an edge is labelled.  If
$z\in I\cap F$ spanned an edge of $L\setminus L'$, that edge must have lost
its last unit in Claim~\ref{clm:transfer}, the only operation that depletes a
labelled edge.  The claim then sends the unit to the edge's unique spanner
$z$, contradicting that $z$ still has charge zero.
\end{proof}

\begin{claim}[internal residual zeros]\label{clm:vonezero}
Every vertex of $I\cap V_1$ spans only edges of $L'$.
\end{claim}

\begin{proof}[Proof of the claim]
Fix $v'\in I\cap V_1$ and let $Y$ be its packing component.  We first show
that every neighbour of $v'$ lies in $V(Y)$.  Partition the incident edges at
$v'$ into those whose other endpoint lies in $V(Y)$ and those whose other
endpoint does not.  The first part is nonempty.  No triangle can use one edge
from each part.  For if $x'v'u'$ were such a triangle and
$x'v'\notin E(S)$, then none of its three edges could belong to $E(S)$---any
such packing edge would join the two packing components---contradicting
maximality.  If $x'v'\in E(S)$, write $x'v'y'\in S$.  The component $Y$
contains only excellent triangles: a zero in a component meeting $S_3$ is
impossible, and a zero in a component meeting $S_2$ lies in $V_2$, whereas
$v'\in V_1$.  Thus $x'v'y'\in S_1$, and its labelled edge is $x'y'$ because
$v'\notin V(L)$.  Let $q'\in F$ span $x'y'$.  If $q'=u'$, then $q'$ also
spans $x'v'$, giving two labelled edges in $x'v'y'$.  If $q'\ne u'$, then
$q'x'y'$ and $u'x'v'$ replace $x'v'y'$ by (E1).  Both alternatives are
impossible.  Split-normality now forces the second part of the incident star
to be empty, proving $N(v')\subseteq V(Y)$.

Suppose that $v'$ spans an edge $ab$.  Then $a,b\in V(Y)$.  If $ab\in L$,
it is the labelled edge of an excellent packing triangle in $Y$ and was not
used by Claim~\ref{clm:transfer}, which uses only edges of $S_2$; hence
$ab\in L'$.  Assume instead that $ab\notin L$.  If $ab\in E(S)$, write
$abd\in S_1$.  When $d=v'$, this packing triangle has no labelled edge,
contrary to its membership in $S_1$.  When $d\ne v'$, the triangle $abv'$ is
in $G-L$, so the unlabelled edge $ab$ of the excellent triangle $abd$ lies in
another triangle of $G-L$, again a contradiction.

It remains that $ab\notin E(S)\cup L$.  Maximality applied to the triangle
$abv'$ forces one of $av'$ and $bv'$ into $E(S)$; say $bv'\in E(S)$ and
write $bv'c\in S_1$.  Here $c\ne a$.  The triangle $abv'$ lies in $G-L$, so
the unlabelled edge $bv'$ of the excellent triangle $bv'c$ again lies in
another triangle of $G-L$.  This final contradiction proves the claim.
\end{proof}

We finish with the crown step, spelling out the matching equality.  Form the
bipartite graph $B$ with left part $I$, right part
\[
  S_G(I)=\{e\in L': e\text{ is spanned by a vertex of }I\},
\]
and incidence $ce\in E(B)$ exactly when $c$ spans $e$.  If $|I|>|L'|$, then
a maximum matching of $B$ leaves a left vertex unmatched.  Start alternating
paths at all unmatched left vertices, traverse unmatched edges left-to-right
and matched edges right-to-left, and let $C$ and $H$ be the reached left and
right vertices.  Every neighbour of a reached left vertex is reached; every
reached right vertex is matched, or the path would augment; and its matched
mate is reached.  Consequently
\[
  N_B(C)=H=S_G(C),
\]
and the matching saturates $H$.  In particular $C$ is nonempty.

For each matched pair $ce$, choose the triangle consisting of $c$, the edge
$e$, and its endpoints.  Notice that $C\cap V(H)=\varnothing$: vertices of
$F$ lie outside $V(S)$, vertices of $V_1$ lie outside $V(L)$, and every edge
of $H\subseteq L$ has both endpoints in $V(L)$.  The chosen triangles are
therefore edge-disjoint: their head edges and crown vertices are distinct,
and $C$ is independent by Claim~\ref{clm:independent}.  Finally, if $c\in C$
had a neighbour outside
$C\cup V(H)$, triangle support of that incident edge would give an edge
spanned by $c$ outside $H=S_G(C)$, a contradiction.  Thus
$(C,H,V(G)\setminus C)$ is an actual fat-head crown, contrary to (c).

Therefore $|I|\leq |L'|$.  Choose an injection from $I$ to $L'$ and move the
unit on each selected edge to its assigned vertex; unused edge charge may
remain.  Every graph vertex now has at least one unit, while total charge
remained $3|S|$.  Hence $|V(G)|\leq3|S|$.
\end{proof}

\begin{proof}[Proof of Proposition~\ref{prop:twelve}]
Assume $\nu(G)=4$ and $\tau(G)\geq9$.  Replace $G$ by its triangle support
and split vertices until the graph is split-normal.  These operations preserve
the triangle hypergraph and therefore preserve $\nu$ and $\tau$; splitting can
only increase the number of vertices.

A triangle-isolated $K_4$ edge block would, by
Lemma~\ref{lem:k4-offset}, leave a graph $G'$ with
\[
  \nu(G')=3,
  \qquad
  \tau(G')\geq7>2\nu(G'),
\]
contradicting Parker's theorem~\eqref{eq:parker}.  Similarly, a nonempty crown
of size $h$ would, by Lemma~\ref{lem:crown-offset}, leave
\[
  \nu(G')=4-h\leq3,
  \qquad
  \tau(G')-2\nu(G')
    =\tau(G)-8+h>0,
\]
again contradicting~\eqref{eq:parker}.

Choose a maximal packing and saturate it under (E1)--(E3).  All hypotheses of
Lemma~\ref{lem:discharge} now hold, so the split graph has at most
$3|S|\leq3\nu(G)=12$ vertices.  The original triangle support has no more
vertices than its split graph, proving the proposition.
\end{proof}

\section{Marked four-packing types}\label{sec:types}

Let $P=\{T_1,T_2,T_3,T_4\}$ be four edge-disjoint triangles.  Any two members
of $P$ share at most one vertex.  For every vertex lying in at least two
members, record the non-singleton subset of $[4]$ consisting of the triangles
that contain it.  No pair $\{i,j\}$ occurs in two recorded subsets, and every
index occurs at most three times.  Private vertices fill the remaining
incidences so that every $T_i$ has size three.

Conversely, any family of non-singleton subsets of $[4]$ with these two
properties defines a linear four-edge, $3$-uniform hypergraph, hence a marked
four-packing type.  There are at most six shared vertices because they consume
at least two of the twelve incidences.  We quotient only by the 24
permutations of the marked triangles.  We do not quotient the ambient graph.

For an ambient order $n$, include every marked type whose union uses at most
$n$ vertices; the remaining vertices are outside the marked packing.  The 16
canonical incidence families are listed in Appendix~\ref{app:types}.
Counting each family at every ambient order at least its support order gives
Table~\ref{tab:types}.

\begin{table}[ht]
\centering
\caption{Marked support types used at each ambient order.}\label{tab:types}
\begin{tabular}{crrrrrrrr}
\toprule
order $n$ & 6 & 7 & 8 & 9 & 10 & 11 & 12 & total \\
\midrule
types & 1 & 3 & 6 & 11 & 14 & 15 & 16 & 66 \\
\bottomrule
\end{tabular}
\end{table}

Orders below six are impossible: four edge-disjoint triangles use twelve
distinct graph edges, while $K_5$ has only ten.  Every graph with $\nu(G)=4$
contains a maximum marked four-packing.  Relabelling its packing blocks and
vertices places it in one of the 66 order-type scopes above.  Thus the marked
scope is complete for every triangle-supported graph of order at most twelve.

\section{The SAT encoding}\label{sec:sat}

Fix an order $n$ and one marked type.  For every possible graph edge
$e\in\binom{[n]}2$, introduce a variable $x_e$.  For every possible triangle
$T\in\binom{[n]}3$, introduce a variable $y_T$.  The base formula contains:
\begin{align}
  y_T &\longrightarrow x_e &&(e\in\tbinom{T}{2}),\label{eq:tri-forward}\\
  \bigwedge_{e\in\binom{T}{2}}x_e &\longrightarrow y_T,
      \label{eq:tri-backward}\\
  x_e &\longrightarrow \bigvee_{T\supset e}y_T,
      \label{eq:edge-support}\\
  \bigvee_{T\ni v}y_T &&& (v\in[n]),\label{eq:vertex-support}
\end{align}
and a unit clause $y_T$ for each of the four marked triangles.
Equations~\eqref{eq:tri-forward}--\eqref{eq:tri-backward} make the $y$-variables
the complete literal triangle closure of the graph selected by the
$x$-variables.  Equations~\eqref{eq:edge-support}--\eqref{eq:vertex-support}
make the graph triangle-supported.

The incremental search uses only two kinds of learned semantic clause.

\begin{lemma}[packing cut]\label{lem:packing-cut}
If a decoded model contains five edge-disjoint triangles
$T_1,\ldots,T_5$, then
\[
  \neg y_{T_1}\vee\cdots\vee\neg y_{T_5}
\]
is valid for every graph with $\nu\leq4$.
\end{lemma}

\begin{lemma}[cover cut]\label{lem:cover-cut}
If a decoded model has a triangle transversal $C$ with $|C|\leq8$, then
\[
  \bigvee_{T:\,E(T)\cap C=\varnothing} y_T
\]
is valid for every graph with $\tau\geq9$.
\end{lemma}

\begin{proof}
The packing clause merely says that the five displayed triangles cannot all
exist.  For the cover clause, a graph with no triangle disjoint from $C$ has
$C\cap E(G)$ as a transversal of size at most eight, even if some edges of
$C$ disappear in a later assignment.
\end{proof}

The algorithm solves the base formula, reconstructs the literal graph and its
complete triangle set, and searches by exact recursion first for a
five-packing and then for a transversal of size at most eight.  It appends the
corresponding valid clause and repeats; no additional graph-specific clause is
used.  When the final formula is unsatisfiable, every
triangle-supported graph in that marked scope has either a five-packing or an
eight-edge transversal.

\section{Certificate verification}\label{sec:computation}

Each case is encoded directly as a SAT instance over the possible graph
edges, with semantic clauses obtained from exact packing and covering
calculations.  The Lean development~\cite{Lean4,Mathlib} reconstructs every
serialized witness,
validates all 331,430 semantic clauses, checks their exact agreement with the
final DIMACS formulas, and discharges each of the 66 formulas by a
kernel-checked LRAT refutation.  An independent executable verifier performs
the same witness reconstruction and formula comparison, and the original
DRAT traces are checked separately with \texttt{drat-trim}.

A separate coverage checker reconstructs every expected marked type at each
order and verifies that the resulting 66 cases occur exactly once.

Table~\ref{tab:campaign} records the complete computation.  Here ``models''
means decoded assignments that produced a semantic cut; it does not denote
isomorphism classes.

\begin{table}[ht]
\centering
\caption{Certified finite computation.}\label{tab:campaign}
\begin{tabular}{rrrrrr}
\toprule
order & cases & models & five-packing cuts & cover cuts & result \\
\midrule
6  & 1  & 4       & 0       & 4      & UNSAT \\
7  & 3  & 220     & 96      & 124    & UNSAT \\
8  & 6  & 2,955   & 1,865   & 1,090  & UNSAT \\
9  & 11 & 18,660  & 12,963  & 5,697  & UNSAT \\
10 & 14 & 50,781  & 36,413  & 14,368 & UNSAT \\
11 & 15 & 95,176  & 67,324  & 27,852 & UNSAT \\
12 & 16 & 163,634 & 113,286 & 50,348 & UNSAT \\
\midrule
total & 66 & 331,430 & 231,947 & 99,483 & UNSAT \\
\bottomrule
\end{tabular}
\end{table}

All 66 DRAT proofs were accepted by \texttt{drat-trim}; together they contain
570,595,138 bytes.  Their 572 binary LRAT segments occupy 250.98 MiB.  Each
formula, proof, computation record, and verification program is identified by
its SHA-256 digest.

The formal proof covers the graph-theoretic definitions, Parker's
packing-rank-three bound, the twelve-vertex reduction, the classification of
marked four-packings, the graph-to-CNF interpretation, and the final case
dispatch.  The SAT solver, witness-generation scripts, and DRAT-to-LRAT
converter are not trusted: their outputs are checked by the formal
development.  The trusted base consists of the Lean kernel, the imported
Mathlib definitions and proof infrastructure, and the checked source and
certificate data.

\begin{proposition}\label{prop:finite}
Every triangle-supported graph $G$ on at most twelve vertices with
$\nu(G)=4$ has $\tau(G)\leq8$.
\end{proposition}

\begin{proof}
Mark a maximum four-packing.  By Section~\ref{sec:types}, the marked graph
belongs to one of the 66 scopes.  The corresponding final formula is
DRAT-refuted.  Lemmas~\ref{lem:packing-cut} and~\ref{lem:cover-cut} show that
every graph with $\nu\leq4$ and $\tau\geq9$ satisfies every learned clause.
Hence no such graph exists in the marked scope.
\end{proof}

\section{Proof of the main theorem}

\begin{proof}[Proof of Theorem~\ref{thm:main}]
If $\nu(G)\leq3$, the result is Parker's theorem~\eqref{eq:parker}.  Suppose
$\nu(G)=4$ and, for a contradiction, $\tau(G)\geq9$.  Put
$H=\supp(G)$.  Proposition~\ref{prop:twelve} and~\eqref{eq:support} give
$|V(H)|\leq12$, $\nu(H)=4$, and $\tau(H)=\tau(G)$.  Applying
Proposition~\ref{prop:finite} to the triangle-supported graph $H$ gives
$\tau(G)=\tau(H)\leq8$, a contradiction.
\end{proof}

\section{Concluding remarks}

Theorem~\ref{thm:main} extends the exact inequality from packing number three
to packing number four.  The reduction in Proposition~\ref{prop:twelve} is
specific to this bounded case and does not provide an induction on the packing
number.  Extending the method beyond rank four will therefore require an
additional structural principle.

\paragraph{Acknowledgments.}
ChatGPT 5.6 Sol (OpenAI) and Claude Fable 5 (Anthropic) assisted with research,
exposition, software, and formalization.  The author checked their output and
takes responsibility for the paper.

\appendix

\section{Canonical marked-packing incidence families}\label{app:types}

For compactness, write $13$ for the shared vertex belonging to marked
triangles $T_1$ and $T_3$, and similarly write $234$ or $1234$ for a vertex
shared by three or four marked triangles.  The complete canonical list is:

\begin{center}
\begin{tabular}{cl}
\toprule
support order & non-singleton incidence subsets \\
\midrule
6  & $12,13,23,14,24,34$ \\
7  & $12,13,23,14,24$ \\
7  & $12,13,14,234$ \\
8  & $12,13,23,14$ \\
8  & $12,13,24,34$ \\
8  & $12,13,234$ \\
9  & $12,13,23$ \\
9  & $12,13,14$ \\
9  & $12,13,24$ \\
9  & $12,134$ \\
9  & $1234$ \\
10 & $12,13$ \\
10 & $12,34$ \\
10 & $123$ \\
11 & $12$ \\
12 & none \\
\bottomrule
\end{tabular}
\end{center}

Every row satisfies the pair-occupancy and degree conditions in
Section~\ref{sec:types}.  Conversely, recursively add a non-singleton subset
of $[4]$ only when it repeats no pair and raises no block degree above three;
then quotient the resulting family by the 24 permutations of $[4]$.  This
produces exactly the rows above.  Private vertices fill each block to degree
three, giving the displayed support order.  Taking all rows of support order
at most $n$ yields the cumulative counts $1,3,6,11,14,15,16$ for
$n=6,\ldots,12$ and therefore 66 order-type scopes.

\small
\bibliographystyle{plain}
\bibliography{tuza-rank-four}

\end{document}
